\documentclass[11pt]{article}

\usepackage[margin=1.15in]{geometry}
\usepackage{amsmath,amssymb,amsthm}
\usepackage{booktabs}
\usepackage{microtype}
\usepackage{xcolor}
\usepackage[colorlinks=true,linkcolor=blue!50!black,citecolor=blue!50!black,urlcolor=blue!50!black]{hyperref}

\newtheorem{theorem}{Theorem}[section]
\newtheorem{lemma}[theorem]{Lemma}
\newtheorem{corollary}[theorem]{Corollary}
\newtheorem{proposition}[theorem]{Proposition}
\newtheorem{remark}[theorem]{Remark}
\newtheorem{question}[theorem]{Question}
\theoremstyle{definition}
\newtheorem{definition}[theorem]{Definition}

\newcommand{\N}{\mathbb{N}}
\newcommand{\Z}{\mathbb{Z}}
\newcommand{\Q}{\mathbb{Q}}
\newcommand{\R}{\mathbb{R}}
\newcommand{\Ztwo}{\Z_2}
\newcommand{\Qtwo}{\Q_2}
\newcommand{\code}[1]{\texttt{#1}}
\newcommand{\ord}{\operatorname{ord}}

\title{Rung One of the Word-Complexity Ladder:\\
  a Machine-Verified Exclusion of an Automatic\\
  Collatz Itinerary by Mahler's Method in Degree One}
\author{M.~Sharpe\thanks{\texttt{msharpe@irendi.com}. The mathematics,
  the Lean formalization, and this text were produced in collaboration
  with Claude (Anthropic). The Collatz conjecture remains open, and
  nothing in this paper claims otherwise.}}
\date{August 28, 2026}

\begin{document}
\maketitle

\begin{abstract}
Let $w = 110\,110\,011\,011\,110\,011\ldots$ be the fixed point of the
substitution $\sigma\colon 1 \mapsto 110$, $0 \mapsto 011$. Its density
of ones is $2/3 > \log 2/\log 3$, so it lies inside the band where a
divergent orbit of the Terras map $T$ is allowed to live. We prove
that no positive integer has Terras parity itinerary $w$, or any shift
of $w$. Equivalently, the $2$-adic Mahler value $F_2(8/9)$, where
$F(z) = \sum_i w_i z^i$ satisfies
$F(z) = (z+z^2)/(1-z^3) + (1-z^2)F(z^3)$, is irrational. The proof is
Mahler's method with a degree-one auxiliary form, carried out entirely
in the language of integer divisibility: a Pad\'e pair
$(p_0, p_1) = (1+z^3-z^4,\ -1+z-z^2+z^3)$ found on the first ten
letters of $w$, the functional equation iterated along
$\alpha_k = (8/9)^{3^k}$, and a height ledger that closes because
$2^{27} > 28\cdot 9^7$. Every step is machine-verified in Lean~4
(\code{Collatz/Rung1.lean}, axiom-free). The earlier reduction paper
recorded the truncation approach to $F_2(8/9)$ as provably
insufficient (irrationality exponent $0.946 < 1$) and the $2$-adic
transposition of Mahler's method as the open work item; we show the
transposition is short --- in the ultrametric setting the analytic
lemma is trivial and non-vanishing is free, so neither a zero estimate
nor the transcendence of $F$ over $\Q(z)$ is needed. This is the first
exclusion of a genuinely aperiodic itinerary class for the $3x+1$
problem, and, to our knowledge, the first formalization of a
Mahler-method argument in a proof assistant. We are explicit about
what is classical: the irrationality of $p$-adic Mahler values at
rational points inside the unit disc is a known type of result.
Finally we begin rung two: an elementary prefix-power criterion
(the Diophantine content of the parity bijection, machine-verified)
excludes every Sturmian itinerary whose slope has partial quotients
$\ge 2$ infinitely often, for slopes up to $0.789$ --- conditional on
two classical facts, Lagrange's best-approximation property and the
three-distance theorem, which are isolated as hypotheses of a
machine-verified level theorem.
\end{abstract}

\section{Introduction}

The word-complexity ladder \cite{ladder} orders potential divergent
orbits of the Terras map
\[
T(n) = \begin{cases} n/2 & n \text{ even},\\ (3n+1)/2 & n \text{ odd}\end{cases}
\]
by the complexity of their parity itineraries $w_t = T^t(n) \bmod 2$.
Rung zero --- a divergent orbit's itinerary is never eventually
periodic --- is formal. The next rung consists of \emph{automatic}
itineraries, fixed points of constant-length substitutions, and
\cite{ladder} reduced the running example exactly to one arithmetic
statement: for $w = \sigma^\infty(1)$ with $\sigma(1) = 110$,
$\sigma(0) = 011$,
\begin{equation}\label{eq:reduction}
\text{some integer has itinerary } w
\iff F_2(8/9) \in \Q,
\end{equation}
where $F(z) = \sum_{i \ge 0} w_i z^i \in \Z[[z]]$ and $F_2$ is
evaluation in $\Q_2$ (convergent since $|8/9|_2 = 1/8$). The
reduction paper also proved that the obvious approach fails: the
partial sums of $F_2(8/9)$ have denominators $9^N$ and $2$-adic error
$2^{-3N}$, an irrationality exponent of $(3\log 2)/(2 \log 3) = 0.946
< 1$, so no Liouville-type argument from truncation can decide
\eqref{eq:reduction}. It identified Mahler's method as the route and
its $2$-adic transposition as ``bounded, classical work, but real
work.''

This paper does that work, and finds it shorter than expected.

\begin{theorem}[Rung one, running example]\label{thm:main}
For every $n \in \N$ and every $s \in \N$, the Terras itinerary of $n$
is not the $s$-shift of $w$: it is false that $T^t(n) \bmod 2 =
w_{s+t}$ for all $t$. In particular no Collatz orbit has itinerary
$w$, and no orbit ever synchronises with $w$.
\end{theorem}

\begin{theorem}[Irrationality]\label{thm:irr}
$F_2(8/9) \notin \Q$.
\end{theorem}

Both are formalized: \code{no\_shifted\_sigWord\_itinerary} and
\code{core} in \code{lean/Collatz/Rung1.lean}, depending only on
\code{propext}, \code{Classical.choice}, \code{Quot.sound}, with no
\code{sorry}.

\paragraph{What is new, and what is not.} We want to be precise.
\begin{itemize}
\item \emph{New as a Collatz statement.} Theorem~\ref{thm:main} is the
first exclusion of an aperiodic itinerary class: rung zero excluded
periodic words, and Theorem~\ref{thm:main} excludes an automatic word of
supercritical density together with all its tails. The density
corridor \cite{companion} permits such an orbit; the arithmetic of
$F$ forbids it.
\item \emph{New as a formalization.} To our knowledge no Mahler-method
argument (transcendence or irrationality of values of Mahler
functions) has previously been checked by a proof assistant. The
formalization is possible because the whole argument is stated with
integers and divisibility by powers of $2$; no $p$-adic analysis and
no Mathlib theory of $\Z_2$ is used.
\item \emph{Not new as number theory.} That values of Mahler functions
at rational (or algebraic) points inside the $p$-adic unit disc are
irrational or transcendental is a classical theme: Mahler
\cite{mahler}, Loxton and van der Poorten \cite{lvdp}, Bundschuh and
V\"a\"an\"anen \cite{bv} and others. Our Section~\ref{sec:classical}
places the proof in that context. What we add is the observation, of
some independent interest, that for \emph{irrationality} at a point
with $|\alpha|_p < 1$ the method needs only a degree-one auxiliary
form, and that the two ingredients which make the archimedean method
delicate --- the analytic lemma and the non-vanishing of the auxiliary
function along the orbit --- are both trivial ultrametrically.
\item \emph{The real content of the $0.946$ barrier.} The truncation
route trades height $9^L$ for $2$-adic gain $2^{3L}$. A degree-one
Pad\'e form trades height $9^{4 \cdot 3^k}$ (plus a sunk cost
$9^{(5/2)3^k}$ independent of the degree) for gain $2^{9 \cdot
3^{k+1}}$: per unit of degree the ratio is $6 \log 2/\log 9 = 1.89 >
1$. This is the classical ``squaring'' that Mahler's iteration buys;
the barrier was a fact about truncation, not about the number.
\end{itemize}

\paragraph{Organisation.} Section~\ref{sec:setup} fixes notation and
proves the exact finite form of the functional equation and of the
block structure of the correction term. Section~\ref{sec:tower} sets
up the tower of forced rationals. Section~\ref{sec:pade} gives the
certificate and the exact $2$-adic size. Section~\ref{sec:ledger}
does the height accounting and proves the theorems.
Section~\ref{sec:classical} discusses the relation to Mahler's method
and the extent of generality. Section~\ref{sec:next} records the next
rung. Throughout, lemma names in \code{typewriter} refer to the Lean
file.

\section{Setup: the word, prefix sums, and the exact identities}
\label{sec:setup}

\subsection{The word}

Define $w\colon \N \to \{0,1\}$ by $w_0 = 1$ and, for $q \ge 0$,
\begin{equation}\label{eq:word}
w_{3q} = w_q, \qquad w_{3q+1} = 1, \qquad w_{3q+2} = 1 - w_q
\end{equation}
(\code{sigWord}, \code{sigWord\_three}, \code{sigWord\_three\_one},
\code{sigWord\_three\_two}). This is the fixed point $\sigma^\infty(1)$:
the block $w_{3q}w_{3q+1}w_{3q+2}$ is $\sigma(w_q)$. The first ten
letters are $1,1,0,1,1,0,0,1,1,1$ (\code{sigWord\_first\_ten}); every
block has exactly two ones.

\subsection{Homogenised prefix sums}

For $a, b \in \Z$ and $L \in \N$ put
\[
S_L(a,b) = \sum_{q < L} w_q\, a^q\, b^{L-1-q}, \qquad
G_L(a,b) = \sum_{q < L} a^q\, b^{L-1-q},
\]
so that $S_L(a,b) = b^{L-1} P_L(a/b)$ where $P_L(z) = \sum_{q<L} w_q
z^q$ is the degree-$<L$ prefix of $F$, and $G_L$ is the same for the
all-ones word. Both are defined by the back-recursions $S_0 = 0$,
$S_{L+1} = b\,S_L + w_L a^L$ and $G_0 = 0$, $G_{L+1} = b\,G_L + a^L$
(\code{preS}, \code{geomS}). The geometric identity
\begin{equation}\label{eq:geom}
(b-a)\,G_L(a,b) = b^L - a^L
\end{equation}
is immediate by induction (\code{geomS\_mul}).

\begin{lemma}[Functional equation, finite form; \code{preS\_three}]
\label{lem:fe}
For all $a, b \in \Z$ and $L \in \N$,
\[
S_{3L}(a,b) = (ab + a^2)\, G_L(a^3, b^3) + (b^2 - a^2)\, S_L(a^3, b^3).
\]
\end{lemma}

\begin{proof}
Induction on $L$. Unrolling the recursion three times,
$S_{3L+3} = b^3 S_{3L} + w_{3L} a^{3L} b^2 + w_{3L+1} a^{3L+1} b +
w_{3L+2} a^{3L+2}$, and by \eqref{eq:word} the three new terms are
$a^{3L}\big(w_L b^2 + ab + (1-w_L)a^2\big) = a^{3L}\big(ab + a^2 +
w_L(b^2-a^2)\big)$. On the right, $G_{L+1}(a^3,b^3) = b^3 G_L + a^{3L}$
and $S_{L+1}(a^3,b^3) = b^3 S_L + w_L a^{3L}$, and the two sides agree
by the induction hypothesis and a ring identity.
\end{proof}

Dividing by $b^{3L-1}$ and letting $L \to \infty$ formally, this is
$F(z) = (z+z^2)/(1-z^3) + (1-z^2)F(z^3)$ at $z = a/b$; but we never
pass to the limit. Every statement below is about finite $L$.

\subsection{The correction term of a word, and the block formula}

For an arbitrary word $u\colon \N \to \{0,1\}$ define
\[
d_u(0) = 0, \qquad d_u(T+1) = 3^{u_T}\, d_u(T) + 2^T u_T, \qquad
j_u(T) = \sum_{t<T} u_t
\]
(\code{wordD}, \code{wordJ}). If the itinerary of $n$ agrees with $u$
on $t < T$, then $d_u(T)$ is the correction term $d(T,n)$ of the exact
orbit identity $2^T\,T^T(n) = 3^{j}\,n + d$ of \cite{companion}, and
$j_u(T)$ is its odd-step count $j$ (\code{dcoef\_eq\_wordD}; this is
the cocycle law of \cite{companion} at $T = 1$). The word-level
cocycle law is
\begin{equation}\label{eq:cocycle}
d_u(s+T) = 3^{j_{u'}(T)}\, d_u(s) + 2^s\, d_{u'}(T), \qquad
j_u(s+T) = j_u(s) + j_{u'}(T), \qquad u' = u(s + \cdot)
\end{equation}
(\code{wordD\_add}).

\begin{lemma}[Block formula; \code{wordD\_sigWord\_three}]
\label{lem:block}
For all $L$, $j_w(3L) = 2L$ and
\[
d_w(3L) = 10\,(9^L - 8^L) - 5\, S_L(8,9).
\]
\end{lemma}

\begin{proof}
Induction on $L$: appending the block $\sigma(w_L)$ to a prefix of
length $3L$ with correction $d$ gives $9d + 5\cdot 8^L$ if $w_L = 1$
(append $1,1,0$) and $9d + 10 \cdot 8^L$ if $w_L = 0$ (append $0,1,1$),
i.e.\ $d_w(3L+3) = 9\,d_w(3L) + (10 - 5w_L)8^L$, which matches the
recursion $S_{L+1}(8,9) = 9 S_L(8,9) + w_L 8^L$.
\end{proof}

\subsection{From an itinerary to a congruence for every prefix}

\begin{proposition}[The link; \code{link}]\label{prop:link}
Suppose the itinerary of $n$ is the $s$-shift of $w$. Put
\[
q = 3^{j_w(s)}, \qquad p = 2^s n - d_w(s), \qquad
D_0 = 5q, \qquad u_0 = 9\,(p + 10q).
\]
Then for every $L$ with $3L \ge s$,
\begin{equation}\label{eq:link}
2^{3L} \ \Big|\ 9^L (p + 10 q) - 5 q\, S_L(8,9).
\end{equation}
\end{proposition}

\begin{proof}
Write $3L = s + T$. By \eqref{eq:cocycle} with $u = w$, the
identification $d_{w'}(T) = d(T,n)$, $j_{w'}(T) = j$ for the orbit of
$n$, and the exact orbit identity $2^T\,T^T(n) = 3^j n + d(T,n)$,
\[
3^{j_w(s)}\, d_w(3L) = 3^{j_w(s) + j}\,\big(d_w(s) - 2^s n\big) +
2^{s+T}\, 3^{j_w(s)}\, T^T(n),
\]
and $j_w(s) + j = j_w(3L) = 2L$. Substituting the block formula for
$d_w(3L)$ and using $2^{3L} \mid 10 \cdot 8^L q$ gives \eqref{eq:link}.
\end{proof}

In $2$-adic language, \eqref{eq:link} for all $L$ says exactly that
$F_2(8/9) = u_0/(9 D_0) \cdot 9 = (p + 10q)/q \cdot 9/5$, i.e.\ that
$x(w) = p/q$ is the rational with odd denominator $3^{j_w(s)}$ obtained
by pulling $n$ back $s$ steps through the $2$-adic Terras map. Shifts
therefore cost nothing: they only change the odd denominator. For
$s = 0$, $q = 1$ and $p = n$.

\section{The tower}\label{sec:tower}

Put $a_k = 8^{3^k} = 2^{3^{k+1}}$ and $b_k = 9^{3^k}$
(\code{lvA}, \code{lvB}), so $\alpha_k = a_k/b_k = (8/9)^{3^k}$,
$a_{k+1} = a_k^3$, $b_{k+1} = b_k^3$, $|\alpha_k|_2 = 2^{-3^{k+1}}$.
Given $D_0, u_0 \in \Z$ define (\code{lvD}, \code{lvU})
\begin{align*}
D_{k+1} &= D_k\,(b_k^2 - a_k^2)(b_k^3 - a_k^3),\\
u_{k+1} &= b_k^2\Big((b_k^3 - a_k^3)\,u_k - D_k\, a_k (a_k + b_k)\, b_k\Big).
\end{align*}
These are the numerator and denominator obtained by running the
functional equation backwards: if $F_2(\alpha_k) = u_k/D_k$ then
$F_2(\alpha_{k+1}) = (F_2(\alpha_k) - A(\alpha_k))/B(\alpha_k) =
u_{k+1}/D_{k+1}$ with $A(z) = (z+z^2)/(1-z^3)$, $B(z) = 1 - z^2$. We
never use this interpretation; the finite statement is the following.

\begin{definition}[\code{Hyp}]
$\mathrm{Hyp}_k(L_0)$ is the statement: for all $L \ge L_0$,
\[
2^{3^{k+1}(L+1)} \ \Big|\ u_k\, b_k^{L} - D_k\, S_{L+1}(a_k, b_k).
\]
\end{definition}

For $k = 0$ and $D_0 = 5q$, $u_0 = 9(p+10q)$ this is \eqref{eq:link}
with $L+1$ in place of $L$, so Proposition~\ref{prop:link} gives
$\mathrm{Hyp}_0(s)$ (the first few $L$ are discarded harmlessly; only
large $L$ matter).

\begin{proposition}[Transport; \code{Hyp\_succ}]\label{prop:transport}
$\mathrm{Hyp}_k(L_0) \Rightarrow \mathrm{Hyp}_{k+1}(L_0)$.
\end{proposition}

\begin{proof}
Fix $L \ge L_0$ and apply $\mathrm{Hyp}_k$ at prefix length $3(L+1)$:
$u_k b_k^{3L+2} - D_k S_{3(L+1)}(a_k,b_k) = 2^{3^{k+1}\cdot 3(L+1)}c$
for some $c \in \Z$. Write $a = a_k$, $b = b_k$. Multiplying by
$b^3 - a^3$ and using Lemma~\ref{lem:fe} and \eqref{eq:geom} (with
$(a^3,b^3)$ and $L+1$),
\[
(b^3 - a^3)\big(u_k b^{3L+2} - D_k S_{3(L+1)}(a,b)\big)
= u_{k+1}\, b_{k+1}^{L} - D_{k+1}\, S_{L+1}(a_{k+1}, b_{k+1})
+ D_k (ab + a^2)\, a^{3(L+1)}.
\]
Both the left side and $a^{3(L+1)} = 2^{3^{k+2}(L+1)}$ are divisible
by $2^{3^{k+2}(L+1)}$, hence so is the middle term.
\end{proof}

Note that $b^3 - a^3$ is odd, so nothing is lost by multiplying
through: the divisibility is of the product, and that is all we need
downstream. By induction, $\mathrm{Hyp}_0(L_0)$ implies
$\mathrm{Hyp}_k(L_0)$ for all $k$ (\code{Hyp\_all}).

\section{The Pad\'e certificate and the exact $2$-adic size}
\label{sec:pade}

Let $p_0(z) = 1 + z^3 - z^4$ and $p_1(z) = -1 + z - z^2 + z^3$, and
their degree-$4$ homogenisations
\[
P_0(a,b) = b^4 + a^3 b - a^4, \qquad
P_1(a,b) = -b^4 + a b^3 - a^2 b^2 + a^3 b
\]
(\code{P0}, \code{P1}). The pair was found by solving the $9$
homogeneous linear conditions ``coefficients of $z^0, \ldots, z^8$ in
$p_0 + p_1 F$ vanish'' in the $10$ unknown coefficients; the solution
is unique up to scale and has integer entries.

\begin{lemma}[The certificate; \code{pade\_base}]\label{lem:pade}
As an identity in $\Z[a,b]$,
\[
b^9 P_0(a,b) + P_1(a,b)\, S_{10}(a,b) = a^9\big(-b^4 + a(b^3 + a^2 b)\big).
\]
\end{lemma}

\begin{proof}
Expand $S_{10}$ with the first ten letters of $w$ and compare
coefficients; this is a finite computation, done in Lean by
\code{simp} and \code{ring}.
\end{proof}

In words: $E(z) = p_0(z) + p_1(z) F(z) = -z^9 + z^{10} + z^{12} +
\cdots$ has order exactly $9$ at $z = 0$ with leading coefficient
$-1$. The lemma is the homogeneous form of this restricted to the
prefix of length $10$; longer prefixes only add terms divisible by
$a^{10}$:

\begin{lemma}[\code{pade\_gen}]\label{lem:padegen}
For every $M \ge 0$ there is $R \in \Z[a,b]$ with
\[
b^{9+M} P_0(a,b) + P_1(a,b)\, S_{10+M}(a,b) = a^9\big(-b^{4+M} + aR\big).
\]
\end{lemma}

\begin{proof}
Induction on $M$ from Lemma~\ref{lem:pade}, using $S_{L+1} = b S_L +
w_L a^L$: the new term $P_1 w_{L} a^{L}$ with $L = 10 + M \ge 10$ is
divisible by $a^{10}$.
\end{proof}

\begin{lemma}[Exact $2$-adic size; \code{level\_dvd}]\label{lem:size}
Let $v \ge 1$, $a = 2^v$, $b$ odd, $D$ odd, $M \ge 0$, and suppose
\[
2^{v(10+M)} \ \Big|\ u\, b^{9+M} - D\, S_{10+M}(a,b).
\]
Then $Z := P_0(a,b)\, D + P_1(a,b)\, u$ is a nonzero integer divisible
by $2^{9v}$.
\end{lemma}

\begin{proof}
Write $u b^{9+M} - D S_{10+M} = 2^{v(10+M)} t$ and take $R$ from
Lemma~\ref{lem:padegen}. Then
\[
b^{9+M} Z = D\big(b^{9+M}P_0 + P_1 S_{10+M}\big) + P_1\big(u b^{9+M} -
D S_{10+M}\big) = 2^{9v}\, O, \qquad
O = D\big(-b^{4+M} + 2^v R\big) + P_1\, 2^{v(1+M)}\, t .
\]
$O$ is odd: $D b^{4+M}$ is odd and the other two terms are even. Hence
$Z \ne 0$ (else $O = 0$), and $2^{9v} \mid b^{9+M} Z$ with $b^{9+M}$
odd gives $2^{9v} \mid Z$.
\end{proof}

This lemma is the whole of ``Mahler's analytic lemma'' and ``the zero
estimate'' in our setting. Archimedeanly one bounds $|E(\alpha_k)|$
from above by a Schwarz-lemma argument and must separately show
$E(\alpha_k) \ne 0$ for infinitely many $k$, which is where the
transcendence of $F$ over $\Q(z)$ enters. Ultrametrically the size
of $E(\alpha_k)$ is \emph{exactly} $|\alpha_k|_2^{\,9}$ as soon as
$|\alpha_k|_2 < |c|_2 = 1$ where $c = -1$ is the leading coefficient,
because all higher coefficients are integers. Nothing else is needed.

\section{The height ledger and the proofs}\label{sec:ledger}

\begin{lemma}[Size of $Z$; \code{abs\_Z\_le}]\label{lem:absZ}
If $0 \le a \le b$ and $|D|, |u| \le m$ then $|P_0(a,b) D + P_1(a,b) u|
\le 7 b^4 m$.
\end{lemma}

\begin{proof}
$|P_0| \le 3b^4$ and $|P_1| \le 4 b^4$ termwise.
\end{proof}

\begin{lemma}[Tower heights; \code{lv\_bounds}]\label{lem:heights}
Let $m_0 = \max(|D_0|, |u_0|)$ and $E_k = 3^k\, 9^{3 \cdot 3^k}\, m_0$.
Then $|D_k|, |u_k| \le E_k$ for all $k$.
\end{lemma}

\begin{proof}
$E_0 = 729\, m_0 \ge m_0$. With $0 \le a_k \le b_k$: $|D_{k+1}| \le
E_k b_k^5$ and $|u_{k+1}| \le b_k^2\big(b_k^3 E_k + E_k \cdot 2 b_k^3\big)
= 3 b_k^5 E_k$, while $E_{k+1} = 3 b_k^6 E_k$ (\code{lvE\_succ}).
\end{proof}

Also $D_k$ is odd for all $k$ if $D_0$ is (\code{lvD\_odd}): $b_k$ is
odd and $a_k$ even.

\begin{lemma}[The numeric inequality; \code{final\_ineq}]\label{lem:num}
If $0 \le m_0 \le k$ then $7\, b_k^4\, E_k < 2^{9 \cdot 3^{k+1}}$.
\end{lemma}

\begin{proof}
The left side is $7 \cdot 3^k \cdot 9^{7 \cdot 3^k} m_0 \le 7 \cdot 3^k
\cdot 9^{7\cdot 3^k} k$. Since $k + 1 \le 3^k$ and $k < 9^k$,
\[
7 \cdot 3^k k < 28 \cdot 3^k 9^k = 28 \cdot 27^k \le 28^{k+1} \le 28^{3^k},
\]
so the left side is $< 28^{3^k}\, 9^{7 \cdot 3^k} = (28 \cdot 9^7)^{3^k}
\le (2^{27})^{3^k} = 2^{27 \cdot 3^k}$, using $28 \cdot 9^7 = 133\,923\,132
\le 134\,217\,728 = 2^{27}$.
\end{proof}

The margin is $0.22\%$ per tower level, which is why the degree-$4$
pair (order $9$) is used rather than the degree-$2$ pair (order $5$):
the ledger for general $m$ is $3(2m+1)\log 2$ against $(m + 5/2)\log
9$, positive for every $m \ge 2$, but the crude constants of
Lemmas~\ref{lem:absZ} and~\ref{lem:heights} eat the $m = 2$ margin.

\begin{theorem}[\code{core}]\label{thm:core}
Let $D_0$ be odd, $u_0 \in \Z$, and $L_0 \in \N$. Then
$\mathrm{Hyp}_0(L_0)$ is false.
\end{theorem}

\begin{proof}
Suppose $\mathrm{Hyp}_0(L_0)$. Let $m_0 = \max(|D_0|,|u_0|)$ and $k =
m_0$. By Proposition~\ref{prop:transport}, $\mathrm{Hyp}_k(L_0)$ holds;
apply it at prefix length $10 + M$ where $9 + M = \max(L_0, 9)$. With
$v = 3^{k+1}$, $a_k = 2^v$, $b_k$ odd and $D_k$ odd,
Lemma~\ref{lem:size} gives $2^{9 \cdot 3^{k+1}} \le |Z_k|$ for $Z_k =
P_0(a_k,b_k) D_k + P_1(a_k,b_k) u_k$. Lemmas~\ref{lem:absZ}
and~\ref{lem:heights} give $|Z_k| \le 7 b_k^4 E_k$, and
Lemma~\ref{lem:num} gives $7 b_k^4 E_k < 2^{9 \cdot 3^{k+1}}$.
Contradiction.
\end{proof}

\begin{proof}[Proof of Theorem~\ref{thm:main}]
If the itinerary of $n$ were the $s$-shift of $w$,
Proposition~\ref{prop:link} would give $\mathrm{Hyp}_0(s)$ with $D_0 =
5 \cdot 3^{j_w(s)}$ odd, contradicting Theorem~\ref{thm:core}.
\end{proof}

\begin{proof}[Proof of Theorem~\ref{thm:irr}]
$F_2(8/9) \in \Z_2$, so if it were rational it would be $u_0/D_0$ with
$D_0$ odd, and then $\mathrm{Hyp}_0(0)$ would hold (the partial sums
$S_{L+1}(8,9)/9^L$ converge to $F_2(8/9)$ with error $O(2^{-3(L+1)})$).
Theorem~\ref{thm:core} forbids it. (The Lean file proves
Theorem~\ref{thm:core} as its central statement and derives
Theorem~\ref{thm:main}; the two-line passage from $F_2(8/9) \in \Q$ to
$\mathrm{Hyp}_0(0)$ is not formalized, since $F_2$ itself is not
defined in the file.)
\end{proof}

\section{Relation to Mahler's method; scope of the argument}
\label{sec:classical}

\paragraph{The classical shape.} Mahler's method \cite{mahler,
nishioka} proves transcendence of $F(\alpha)$ for $F$ satisfying
$F(z) = A(z) + B(z)F(z^d)$ by constructing an auxiliary polynomial
$P(z, F(z))$ of degree $m$ in $F$ with high vanishing order $N \sim
m^2$ at $0$ (Siegel's lemma), evaluating along $\alpha^{d^k}$, and
comparing the analytic upper bound $|P(\alpha_k, F(\alpha_k))| \ll
|\alpha_k|^N$ with a Liouville lower bound from the algebraic height,
which grows like $d^k \cdot m$. The two delicate points are a zero
estimate guaranteeing $P(\alpha_k, F(\alpha_k)) \ne 0$ for infinitely
many $k$ (this uses that $F$ is transcendental over $\Q(z)$), and the
Schwarz-type analytic lemma.

\paragraph{What changes $2$-adically, and why degree one suffices.}
For irrationality one takes $m = 1$: $E = p_0 + p_1 F$ with $\deg p_i
\le m'$, $N = 2m' + 1$. Archimedeanly $m = 1$ is hopeless (it recovers
only Liouville-grade approximations). Ultrametrically it is enough
because the height grows like $d^k \cdot m'$ while the gain grows like
$d^k \cdot N \cdot v_2(\alpha)$, and for $\alpha = 8/9$, $d = 3$:
$3(2m'+1)\log 2 > (m' + 5/2)\log 9$ for all $m' \ge 2$. The $5/2$ is
the cost of transporting the rational value down the tower; it is
independent of $m'$, which is exactly why the truncation approach
(which pays the analogous cost at every step) sees exponent $0.946$
while the Pad\'e form sees $1.89$ per unit degree. And, as
Lemma~\ref{lem:size} shows, both delicate points disappear: with
integer coefficients and $|\alpha_k|_2 \to 0$, the value of $E$ has
\emph{exactly} the $2$-adic size of its leading term, so non-vanishing
is automatic and the transcendence of $F$ over $\Q(z)$ is never used.
(It is true --- $w$ is aperiodic --- but the certificate $c = -1 \ne 0$
is the only fact about the leading behaviour of $E$ that the proof
consumes.)

\paragraph{Where this sits in the literature.} Irrationality and
transcendence of $p$-adic values of Mahler functions at points inside
the unit disc, and irrationality measures for them, are classical:
see Loxton--van der Poorten \cite{lvdp}, Bundschuh and V\"a\"an\"anen
\cite{bv}, and the function-field development by Fernandes
\cite{fernandes}. We have not located a statement that covers
$F_2(8/9)$ for this specific $F$ verbatim, but we would be surprised
if the general form --- $F \in \Z[[z]]$ Mahler of degree $d$ over
$\Q(z)$, not rational, $\alpha \in \Q$ with $0 < |\alpha|_p < 1$,
$\Rightarrow F_p(\alpha) \notin \Q$ --- were not known to experts,
and we make no priority claim for it. What we claim is the Collatz
consequence, the explicit finite certificate, and the formalization.

\paragraph{The whole rung.} The argument uses only: (i) $w$ is the
fixed point of a constant-length substitution of length $\ell$ with
every block containing the same number $a$ of ones, which gives the
block formula with $(2^\ell, 3^a)$ in place of $(8, 9)$ and a
functional equation of degree $\ell$ with polynomial coefficients;
(ii) $|2^\ell/3^a|_2 = 2^{-\ell} < 1$; (iii) a Pad\'e certificate,
found by linear algebra on the first $2m'+2$ letters, with odd leading
coefficient (or leading coefficient of known $2$-adic size). The
supercritical condition $2^\ell < 3^a$ is what makes the real value
also converge, but it plays no role in the $2$-adic proof: the
argument excludes \emph{sub}critical uniform-weight automatic words as
itineraries too, which is consistent, since those are already excluded
by the density corridor. A parametric Lean theorem with the
certificate as an explicit input is the natural next formalization,
and would turn rung one from one word into the class.

\section{The next rung: Sturmian itineraries}\label{sec:next}

Rung two consists of morphic and Sturmian itineraries. The Mahler route
continues to non-uniform block weights via Hermite--Pad\'e forms for
vector Mahler systems, with the same trivial ultrametric analytic
lemma. But for Sturmian words there is an elementary route with no
functional equation at all, whose threshold lands exactly on the
critical line; it is machine-verified up to two classical facts of
Diophantine approximation, which we isolate as hypotheses.

\subsection{The prefix-power criterion}

\begin{proposition}[\code{prefix\_power}, \code{prefix\_power\_bound}]
\label{prop:pp}
Let $\ell \le M$. If the itinerary of $n$ is $\ell$-periodic on its
first $M$ letters, i.e.\ $T^t(n) \equiv T^{t+\ell}(n) \pmod 2$ for all
$t + \ell < M$, then either $T^\ell(n) = n$, or
\[
2^{M-\ell} \le \max\big(n, T^\ell(n)\big), \qquad\text{hence}\qquad
2^M \le 2^\ell n + 3^{o}\,(n + 2^\ell),
\]
where $o$ is the number of odd steps among the first $\ell$.
\end{proposition}

\begin{proof}
$n$ and $T^\ell(n)$ have the same first $M - \ell$ parities, so they
are congruent modulo $2^{M-\ell}$ by the parity bijection
\cite{companion}; two distinct naturals congruent modulo $2^{M-\ell}$
differ by at least $2^{M-\ell}$. The second form uses the master
inequality $2^\ell (T^\ell(n) + 1) \le 3^o (n + 2^{\ell - o})$.
\end{proof}

In logarithmic form: if the itinerary begins with an $e$-th power (not
necessarily integral) of a word of length $\ell$ and ones-density
$\delta$, then $(e - 1)\ell \le \delta \ell \log_2 3 + \log_2(n+1) +
O(1)$. The threshold exponent $1 + \delta \log_2 3$ equals $2$
precisely at $\delta = \log 2/\log 3$. Below the critical line squares
are impossible for large $\ell$; above it one needs $e > 2$. This is
the Diophantine content of the parity bijection: the $2$-adic
realizations of periodic words are rationals $d(u)/(2^\ell - 3^o)$,
and an integer cannot lie $2$-adically within $2^{-M}$ of one unless
it is as large as the approximation is good.

We note in passing what this criterion is \emph{not}: for a rational
(here, integer) target, Roth-, Ridout- and Subspace-theorem
improvements of Liouville's inequality have no content, since the
lower bound $|Nq - p| \ge 1$ is already exact. The lever is the word
exponent, not the Diophantine exponent.

\subsection{Sturmian words and their long periodic stretches}

Fix $\alpha \in (0,1)$ irrational and $\rho \in \R$, and let
$w_n = \lfloor (n+1)\alpha + \rho\rfloor - \lfloor n\alpha + \rho\rfloor
\in \{0,1\}$ be the mechanical word; every Sturmian word is of this
form, its ones-density is $\alpha$, and every shift of it is again
mechanical with the same slope. Let $p_k/q_k$ be the convergents of
$\alpha$, with partial quotients $a_k$, and $\delta_k = |q_k\alpha -
p_k|$; the sign of $q_k\alpha - p_k$ alternates, and we call $k$
\emph{positive} when it is $> 0$.

Fix a positive $k$, write $q = q_k$, $\delta = \delta_k$, $Q =
q_{k+1}$, and call $n$ a \emph{visit} if $\{n\alpha + \rho\} \ge 1 -
\delta$. Two elementary facts (\code{floor\_shift}, \code{mech\_shift},
\code{visit\_sep}):
\begin{enumerate}
\item[(i)] $\lfloor (n+q)\alpha + \rho \rfloor = \lfloor n\alpha +
\rho\rfloor + p + [n \text{ is a visit}]$, hence $w_{n+q} = w_n$
unless $n$ or $n+1$ is a visit.
\item[(ii)] If $\|m\alpha\| \ge \delta$ for all $0 < m < Q$, two
visits are at least $Q$ apart: their fractional parts lie in a
half-open interval of length $\delta$, so their difference $n' - n$
satisfies $\|(n'-n)\alpha\| < \delta$.
\end{enumerate}
The hypothesis of (ii) is Lagrange's best-approximation property of
convergents: $q_{k+1}$ is the least $m > 0$ with $\|m\alpha\| <
\|q_k\alpha\|$. The second classical input is the three-distance
theorem \cite{sos}: among any $q_k + q_{k+1}$ consecutive points
$\{n\alpha + \rho\}$, every gap is at most $\|q_k\alpha\|$, so every
window of length $G = q_k + q_{k+1}$ contains a visit.

Consequently, if $n_0$ is the first visit, $n_0 < G$, and the word is
$q$-periodic on the run $[n_0 + 1,\, n_0 + Q - 1)$: a run of length
$Q - 2$ starting at $s = n_0 + 1 \le G$.

\subsection{The size of a tail}

Suppose $N$ has itinerary $w$ with $\alpha$ above the critical line,
i.e.\ $2 \le 3^\alpha$. The exact orbit identity gives $2^s T^s(N) =
3^{j_s} N + d_s$ with $j_s = \lfloor s\alpha + \rho\rfloor - \lfloor
\rho\rfloor \le s\alpha + 1$ and, by the closed form of
$d_s = \sum_{i<s} w_i 2^i 3^{\#\{\text{ones in } (i, s)\}}$
\cite{ladder}, each term at most $2^i \cdot 3^{(s-i-1)\alpha + 1} \le
3^{i\alpha} 3^{(s-i-1)\alpha+1} \le 3 \cdot 3^{s\alpha}$ --- the step
$2^i \le 3^{i\alpha}$ is exactly supercriticality. So
(\code{dcoef\_le\_of\_itin}, \code{tail\_le\_of\_itin})
\[
d_s \le 3s\cdot 3^{s\alpha}, \qquad
2^s\, T^s(N) \le 3^{s\alpha}\,(3N + 3s).
\]

\subsection{The level theorem}

\begin{theorem}[\code{sturmian\_level}]\label{thm:level}
Let $2 \le 3^\alpha$, $q\alpha = p + \delta$ with $0 < \delta < 1$,
$Q \ge 3$, and suppose
\begin{itemize}
\item[(L)] $\|m\alpha\| \ge \delta$ for all $0 < m < Q$;
\item[(T)] every window of $G$ consecutive integers contains a visit;
\item[(S)] for every $s \le G$: $\ 2^{Q - 3 + s} > 3^{(q+s)\alpha + 1}\,(3N + 3s + 1)$.
\end{itemize}
Then $N$ does not have itinerary $w$.
\end{theorem}

\begin{proof}
Let $n_0 < G$ be the first visit, $s = n_0 + 1$, $m = T^s(N)$. By (i)
and (ii) the itinerary of $m$ is $q$-periodic on its first $q + Q - 2$
letters. The orbit of $m$ is not $q$-periodic: otherwise $w$ would be
$q$-periodic from $s$ on, so by (i) the visit indicator would be
constant from $s$ on, contradicting (T) (there is a visit) together
with (ii) (two adjacent visits are impossible as $Q \ge 3$).
Proposition~\ref{prop:pp} gives $2^{q+Q-2} \le 2^q m + 3^o(m + 2^q)$
with $o \le q\alpha + 1$, hence $2^{q+Q-2} \le 2^{q+1} 3^{q\alpha+1}
(m+1)$. Multiplying by $2^s$ and using the tail bound together with
$2^s \le 3^{s\alpha}$,
\[
2^{q+1} 2^{Q-3+s} \le 2^{q+1} 3^{q\alpha + 1} 3^{s\alpha} (3N + 3s + 1),
\]
contradicting (S).
\end{proof}

\begin{corollary}[Sturmian rung, conditional on Lagrange and three-distance]
\label{cor:sturmian}
Let $\alpha \in (\log 2/\log 3, 1)$ be irrational, $\eta = \alpha
\log_2 3 - 1 \in (0, 0.585)$, and suppose that for infinitely many
positive $k$
\[
a_{k+1} \;>\; \frac{1 + 2\eta}{1 - \eta}.
\]
Then no natural number has a Sturmian itinerary of slope $\alpha$,
with any intercept.
\end{corollary}

\begin{proof}
Apply Theorem~\ref{thm:level} at such a level with $q = q_k$, $Q =
q_{k+1}$, $G = q_k + q_{k+1}$; (L) is Lagrange's property, (T) is the
three-distance theorem. For (S), take logarithms: since the coefficient
$\alpha \log_2 3 = 1 + \eta$ of $s$ on the right exceeds $1$, the
worst case is $s = G$, and
\[
(Q - 3 + s) - \big((q+s)\alpha + 1\big)\log_2 3 - \log_2(3N + 3s + 1)
\;\ge\; q_{k+1}(1 - \eta) - q_k(1 + 2\eta) - \log_2\!\big(3N + 3(q_k + q_{k+1}) + 1\big) - 5,
\]
which, as $q_{k+1} \ge a_{k+1} q_k$, is at least $q_k\big(a_{k+1}(1 -
\eta) - (1 + 2\eta)\big) - O(\log q_k + \log N)$ and tends to
$+\infty$ along the chosen levels.
\end{proof}

Concretely, partial quotients $\ge 2$ infinitely often suffice for
$\alpha < 0.789$, $\ge 3$ for $\alpha < 0.883$, $\ge 4$ for $\alpha <
0.946$, and $\ge 6$ for every $\alpha < 1$ --- a set of slopes of full
Lebesgue measure in each range. Only the sign restriction (the level
must be positive) and the golden-type slopes, whose partial quotients
are eventually $1$, are outside the corollary. Numerically
(\code{analysis/sturmian\_prefix\_power.py}) the golden-type slopes are
excluded just as well up to $\alpha \approx 0.9$: the returns to the
visit interval take the two values $q_{k+1}$ and $q_{k+1} + q_k$, the
long return gives a run of length $\ge q_{k+1} + q_k - 2$, and the
margin per $q_k$ is then $\approx (a_{k+1} + 1)(1 - \eta) - (1 +
2\eta) - \eta\,(\text{position of the first long return})/q_k$,
which stays positive for $\eta$ up to about $0.4$. Proving that a long
return occurs early enough is a refinement of the three-distance
theorem that we have not carried out.

\subsection{Scope and honesty}

Theorem~\ref{thm:level} is machine-verified with (L), (T), (S) as
hypotheses. Everything about orbits, itineraries and sizes is inside
Lean; the two Diophantine inputs and the asymptotic passage from (S)
to a condition on partial quotients are not. Both inputs are
nineteenth- and mid-twentieth-century classics with elementary
proofs, and formalizing them (Lagrange's property is close to
Mathlib's existing continued-fraction library; the three-distance
theorem is not in Mathlib) is the natural next formalization target.
With them, Corollary~\ref{cor:sturmian} becomes a theorem of the
repository, and the Sturmian part of rung two closes for the stated
slopes. The morphic, non-Sturmian words of rung two remain for the
Mahler route.

\section*{Reproducibility}

\code{lean/Collatz/Rung1.lean} builds against the repository's pinned
Mathlib with \code{lake build Collatz.Rung1}; \code{\#print axioms}
reports \code{propext}, \code{Classical.choice}, \code{Quot.sound} for
the theorems \code{no\_shifted\_sigWord\_itinerary},
\code{no\_sigWord\_itinerary} and \code{core}. \code{lean/Collatz/PrefixPower.lean} and \code{lean/Collatz/Sturmian.lean}
build the same way and report the same axioms.
\code{analysis/pade\_rung1.py} recomputes the Pad\'e
certificates for degrees $2, 3, 4, 6$, checks the predicted $2$-adic
valuations of $E(\alpha_k)$ against the series to $1053$ bits, and
prints the height ledger.

\begin{thebibliography}{10}

\bibitem{bv} P.~Bundschuh and K.~V\"a\"an\"anen, Arithmetical
properties of the solutions of certain functional equations, and
related papers on $p$-adic Mahler values, 1990s--2010s.

\bibitem{fernandes} G.~Fernandes, M\'ethode de Mahler en
caract\'eristique non nulle, \emph{Ann. Inst. Fourier} \textbf{68}
(2018).

\bibitem{ladder} M.~Sharpe, The word-complexity ladder for Collatz
divergence: a machine-verified rung zero and a Mahler-value reduction
of the automatic rung, companion paper and repository, 2026.

\bibitem{lvdp} J.~H.~Loxton and A.~J.~van der Poorten, Arithmetic
properties of the solutions of a class of functional equations,
\emph{J. Reine Angew. Math.} \textbf{330} (1982), 159--172.

\bibitem{mahler} K.~Mahler, Arithmetische Eigenschaften der
L\"osungen einer Klasse von Funktionalgleichungen, \emph{Math. Ann.}
\textbf{101} (1929), 342--366.

\bibitem{nishioka} Ku.~Nishioka, \emph{Mahler Functions and
Transcendence}, Lecture Notes in Mathematics 1631, Springer, 1996.

\bibitem{sos} V.~T.~S\'os, On the distribution mod 1 of the sequence
$n\alpha$, \emph{Ann. Univ. Sci. Budapest. E\"otv\"os Sect. Math.}
\textbf{1} (1958), 127--134; see also N.~B.~Slater, Gaps and steps
for the sequence $n\theta \bmod 1$, \emph{Proc. Cambridge Philos.
Soc.} \textbf{63} (1967), 1115--1123.

\bibitem{companion} M.~Sharpe, The Collatz conjecture at the critical
line: machine-verified no-go, density, and cycle theorems; the
$3$-adic shadow, companion papers and repository, 2026.

\end{thebibliography}

\end{document}
